% Section 10.1, from lectures/L10/appendix/intro_framing.md.

\section{Frames and serialisation}
\label{c10:sec:framing}

\Chapref{c1:ch} to \ref{c8:ch} ended with state machines controlling a single board: a blinking LED,
a UART on an FPGA. The project starts from a different problem: getting independent pieces of
hardware to exchange information reliably over a shared wire, where neither side can see inside the
other.

This is central to embedded systems because the communication happens as raw byte streams, RAM and
flash are limited, and robustness is critical. Nodes usually communicate over serial interfaces such
as UART, SPI, I2C or RS-485. Data is sent as a sequence of bytes, and interpreting that byte stream
requires a structure. That structure is called a \textbf{frame}.

A frame is a structured representation of data that:
\begin{itemize}
  \item Contains clearly defined fields.
  \item Has a defined length.
  \item Can be validated (for example via a checksum).
  \item Can be serialised into a byte array (\code{std::uint8_t[]}).
  \item Can be deserialised from a byte array.
\end{itemize}

Consider node A sending the value \code{0x2A} to node B over a shared wire, with nothing else agreed
in advance:
\begin{itemize}
  \item How does B know when a new value begins, in a stream of bits that looks the same in the
    middle of a value as it does between two values?
  \item How does B know how long (how many bytes) the value is?
  \item How does B know that the bits it sampled were not corrupted by electrical noise?
  \item What happens if two nodes transmit at the same time?
\end{itemize}

Every real protocol answers these with some combination of the following, each named here after the
field this book's frame carries it in:
\begin{itemize}
  \item \textbf{Framing (\code{SOF})}: a recognisable start marker (and often an end marker), so
    that a receiver can decide where a message begins and where it ends.
  \item \textbf{A length field (\code{LEN})}: so that the receiver knows how much payload to expect.
  \item \textbf{A checksum (\code{CHK})}: extra bits, computed from the payload, letting the
    receiver detect (not necessarily correct) many kinds of transmission error.
  \item \textbf{A media access rule}: what happens when more than one node wants to transmit at the
    same time.
\end{itemize}

A \code{TYPE} field joins them to say what the message means, so that the receiver can act on it.

\subsection*{A small but realistic frame}
To make it concrete: design a frame for sending short messages between nodes on a shared,
byte-oriented link (UART, RS-485, \dots):

\lecfig[0.68]{lectures/L10/appendix/images/example_frame_format.png}{The frame format the rest of
this section builds, field by field.}{c10:fig:frame}

Three rules matter as soon as you start working these fields out by hand:
\begin{itemize}
  \item \textbf{Byte order}: every field wider than a byte (SOF, SEQ, CHK) is big-endian, most
    significant byte first.
  \item \textbf{LEN} counts only the DATA field, never the header and never CHK.
  \item \textbf{CHK} is a 16-bit sum of every byte from SOF through the last payload byte, wrapping
    on overflow. It is computed over the \emph{serialised} byte stream rather than over the field
    values as integers, so SOF contributes \code{0xA5 + 0xF7}, not \code{0xA5F7} treated as a single
    16-bit number.
\end{itemize}

\paragraph{Why DST, SRC and SEQ in particular}
SOF, LEN, TYPE and CHK already answer the four questions above. These three exist because a shared
link with more than two participants raises new ones:
\begin{itemize}
  \item \textbf{DST and SRC}: who the frame is for and who it is from, so that a node can ignore
    traffic addressed elsewhere and a reply can find its way back. CAN answers this entirely
    differently, as \secref{c10:sec:broadcast} shows.
  \item \textbf{SEQ}: matches a reply back to \emph{that particular} request, not some other one in
    the air.
\end{itemize}

\subsection*{Choosing SOF}
\begin{itemize}
  \item A single sync byte is easy to hit by accident: any payload byte that happens to equal it
    looks like the beginning of a frame.
  \item Two bytes are therefore common, and this protocol uses \code{0xA5F7}.
  \item SOF is the \emph{first} thing a parser checks when reading a frame. If it does not match the
    data is discarded, and the parser goes on looking for the next SOF.
  \item Note what that gives and does not give: it makes a false synchronisation much less likely,
    but not impossible, since payload bytes can still happen to spell \code{A5 F7}. The remaining
    risk is why LEN and CHK still have to be checked afterwards.
\end{itemize}

\subsection*{Worked example: PING and PONG}
Suppose node \code{0x01} pings node \code{0x02}, with sequence number \code{0x0020} and no payload:
\begin{itemize}
  \item \code{SOF = 0xA5F7}, \code{LEN = 0x00}, \code{TYPE = 0x00} (Ping), \code{DST = 0x02},
    \code{SRC = 0x01}, \code{SEQ = 0x0020}.
  \item \code{CHK = 0xA5 + 0xF7 + 0x00 + 0x00 + 0x02 + 0x01 + 0x00 + 0x20 = 0x01BF}.
\end{itemize}

\begin{textdrawing}
SOF      LEN   TYPE   DST   SRC   SEQ      DATA   CHK
A5 F7    00    00     02    01    00 20    --     01 BF
\end{textdrawing}

Node \code{0x02} replies with a PONG:
\begin{itemize}
  \item \code{SOF = 0xA5F7}, \code{LEN = 0x00}, \code{TYPE = 0x01} (Pong), \code{DST = 0x01},
    \code{SRC = 0x02}, \code{SEQ = 0x0020}.
  \item \code{CHK = 0xA5 + 0xF7 + 0x00 + 0x01 + 0x01 + 0x02 + 0x00 + 0x20 = 0x01C0}.
  \item Note that \code{DST} and \code{SRC} are swapped, \code{SEQ} is unchanged, and only
    \code{TYPE} and the two addresses differ from the PING: identical arithmetic with different
    inputs.
\end{itemize}

That covers three of the four opening questions, plus the extra ones a shared link with several
nodes raises. The fourth is still open: what happens if two nodes transmit at the same time? CAN
answers it differently, and that is the subject of \secref{c10:sec:can}.

\subsection*{Frame types}
TYPE is one byte on the wire, so a frame type travels as an unsigned integer. Every type is assigned
an id of its own by convention: \code{Ping = 0}, \code{Pong = 1}, \code{StatusReq = 2},
\code{StatusResp = 3}. Further types the protocol acquires simply continue the sequence.

Two consequences worth noting before any code is written:
\begin{itemize}
  \item The received byte is untrusted input. Every value the receiver does not recognise must be
    rejected rather than acted on, exactly like a wrong SOF or a wrong checksum.
  \item Reserving an ID immediately above the last real type, an ``unknown'' value, makes that check
    a single comparison. Without it the validation needs a list of valid IDs that has to be kept in
    step every time a type is added.
\end{itemize}

\subsection*{A simple 16-bit checksum}
For this book's purposes a checksum only has to be simple to implement, simple to test, and good
enough to catch many random bit errors. A 16-bit sum meets all three: start at \code{0}, and for
every byte in the range, \code{sum += byte}. The result is the checksum, modulo $2^{16}$.

\begin{cppcode}
/**
 * @brief Compute a simple 16-bit checksum (sum of bytes modulo 2^16).
 *
 * @param[in] data Pointer to the bytes to checksum.
 * @param[in] dataLen Number of bytes to checksum.
 *
 * @return The 16-bit checksum.
 */
std::uint16_t checksum16(const std::uint8_t* data, const std::size_t dataLen) noexcept
{
    if ((nullptr == data) || (0U == dataLen)) { return 0U; }
    std::uint16_t checksum{};
    for (std::size_t i{}; i < dataLen; ++i) { checksum += data[i]; }
    return checksum;
}
\end{cppcode}

\begin{warning}
\textbf{A checksum's limitations, worth being explicit about.} A checksum detects only errors that
change the arithmetic sum. A single flipped bit always changes the sum, so single-bit errors are
always caught. The real weaknesses lie elsewhere: two whole bytes swapping places leave the sum
unchanged, and \emph{some} combinations of several flipped bits cancel each other out and slip
through undetected. CRCs catch a much broader, mathematically well-understood class of errors, and
are therefore used by CAN instead of a simple sum. They are introduced in
\secref{c10:sec:frameformat} and implemented in \chapref{c13:ch}.
\end{warning}

This is not as strong as a CRC, but it is enough for practising framing, parsing and robust error
handling, which is what this introduction is for.

\subsection*{Serialisation and deserialisation}
\textbf{Serialisation} turns structured data into a sequence of bytes (\code{std::uint8_t[]}) so
that it can be sent over an interface. \textbf{Deserialisation} interprets a received byte array and
extracts the structured data again: the fields are validated first (SOF, length, type, checksum),
and only once every check has passed is the data extracted.

Validating before extracting matters because UART and RS-485 deliver bytes as a stream, so a
receiver routinely sees rubbish between frames, dropped bytes and corrupted frames.

In embedded systems this is nearly always done with C arrays (\code{std::uint8_t buffer[Size]}),
pointers to those arrays at call boundaries (\code{std::uint8_t*}), and an explicit length
(\code{std::size_t}).

Both directions need the same byte order handling on every multi-byte field, which is repetitive and
easy to get subtly wrong. Two function templates remove the repetition. Both constrain \code{T} with
\code{std::is_unsigned}, so wherever they end up they need \code{<type_traits>}:

\begin{cppcode}
#include <type_traits>

/**
 * @brief Write a value to a buffer, most significant byte first.
 *
 *        The caller is responsible for the buffer being large enough.
 *
 * @tparam T The value type. Must be unsigned.
 *
 * @param[out] buf The buffer to write to.
 * @param[in] offset Value offset, i.e. the index at which to write the value.
 * @param[in] value  The value to write.
 */
template <typename T>
void writeToBuf(std::uint8_t* buf, const std::size_t offset, const T value) noexcept
{
    static_assert(std::is_unsigned<T>::value, "Failed to write to buffer: T must be unsigned!");
    constexpr std::size_t len{sizeof(T)};
    constexpr std::size_t msb{len - 1U};
    constexpr std::size_t bitsPerByte{8U};

    for (std::size_t i{}; i < len; ++i)
    {
        const std::size_t shift{bitsPerByte * (msb - i)};
        buf[offset + i] = static_cast<std::uint8_t>((value >> shift));
    }
}

/**
 * @brief Read a value from a buffer, most significant byte first.
 *
 * @tparam T The value type. Must be unsigned.
 *
 * @param[in] buf Buffer to read from.
 * @param[in] offset Value offset in the buffer.
 *
 * @return The retrieved value.
 */
template <typename T = std::uint8_t>
[[nodiscard]] T readFromBuf(const std::uint8_t* buf, const std::size_t offset) noexcept
{
    static_assert(std::is_unsigned<T>::value, "Failed to read from buffer: T must be unsigned!");
    constexpr std::size_t len{sizeof(T)};
    constexpr std::size_t msb{len - 1U};
    constexpr std::size_t bitsPerByte{8U};
    T value{};

    for (std::size_t i{}; i < len; ++i)
    {
        const std::size_t shift{bitsPerByte * (msb - i)};
        value |= static_cast<T>(buf[offset + i]) << shift;
    }
    return value;
}
\end{cppcode}

\begin{warning}
\textbf{A pitfall worth naming, since the compiler does not catch it.} The number of bytes written is
\code{sizeof(T)}, derived from the \emph{argument}, not from the field. Passing in a
\code{std::size_t} where the field is one byte writes \textbf{eight} bytes and silently overwrites
the fields that follow:

\begin{cppcode}
constexpr std::size_t payloadLen{0U};
writeToBuf(buf, FrameOffset::Len, payloadLen);                            // Writes 8 bytes (error).
writeToBuf(buf, FrameOffset::Len, static_cast<std::uint8_t>(payloadLen)); // Writes 1 byte.
\end{cppcode}

Always cast to the field's own width at the call site.
\end{warning}

\subsection*{From paper to code}
This exact frame is made runnable in C++ under \file{lectures/L10/appendix/code/}: a \code{Frame}
struct serialising itself into a byte buffer, deserialising itself back, and rejecting anything with
a wrong SOF, length, type or checksum. Serialising the PING above produces exactly the bytes worked
out above, \code{A5 F7 00 00 02 01 00 20 01 BF}. \file{comm/def.hpp} already provides the field
offsets and field lengths, so you never have to count byte positions by hand.

Turning the type IDs above into a C++ enumeration, declaring the struct and implementing both methods
yourself is Exercise~\ref{c10:ex:cpp}. A set of unit tests comes with the code, in
\file{test/frame_test.cpp}: they check both methods against the frames worked out on paper, byte by
byte. You do not write them; \code{make} runs them, and your job is to make them pass.

\subsection*{What comes next}
One question from the beginning of this section is still open: \textbf{what happens if two nodes
transmit at the same time?}

Nothing in this frame format answers it. A sync word cannot help, since both transmitters send one;
a checksum cannot help, since it only reports the damage after the fact. Answering it takes a
different kind of protocol, one where the bus's \emph{electrical} behaviour decides who continues
and who drops out, before any data is lost. That protocol is CAN.
